\documentclass[../../../include/open-logic-section]{subfiles}

\begin{document}
	
\olfileid[fr]{sth}{z}{nat}
\olsection{Choisir nos nombres naturels}

%Admettons que le principe informel \stagesinf{} puisse être justifié.
%Il vaut cependant aussi la peine de commenter l'axiome particulier
%que nous avons tiré de ce principe informel.

Dans la \olref[infinity-again]{defnomega}, nous avons défini
explicitement l'expression «~nombres naturels~». Comment comprendre
cette stipulation ? Il ne s'agit pas d'une thèse métaphysique, mais
seulement d'une décision de \emph{traiter} certains ensembles comme
les nombres naturels. Nous avons évoqué des raisons de penser ainsi
dans la \olref[sfr][rel][ref]{sec}, la \olref[sfr][arith][ref]{sec} et
la \olref[sfr][infinite][dedekindsproof]{sec}. Mais nous pouvons
formuler ces raisons plus nettement encore.

Notre axiome de l'infini suit \citet{VonNeumann1925}. Nous aurions
toutefois pu adopter, à sa place, l'axiome suivant :

\begin{defish}
\emph{Axiome de l'infini de Zermelo, \citeyear{Zermelo1908Untersuchungen}.}
Il existe un ensemble $A$ tel que $\emptyset \in A$ et
$(\forall x \in A)\{x\} \in A$.
\end{defish}

Si nous avions adopté l'axiome de Zermelo plutôt que le nôtre,
inspiré de von Neumann, nous aurions tout autant obtenu un ensemble
infini au sens de Dedekind, donc une algèbre de Dedekind. Dans
l'approche de Zermelo, l'élément distingué de notre algèbre aurait
encore été $\emptyset$ (notre substitut de $0$), mais l'injection
aurait été donnée par l'application $x \mapsto \{x\}$, au lieu de
$x \mapsto x \cup \{x\}$. La conséquence la plus simple est que
Zermelo prend pour $2$ l'ensemble $\{\{\emptyset\}\}$, tandis que
nous prenons, avec von Neumann, pour $2$ l'ensemble
$\{\emptyset, \{\emptyset\}\}$.

Pourquoi choisir l'un de ces axiomes de l'infini plutôt que l'autre ?
La principale raison pratique est que l'approche de von Neumann
s'étend assez bien aux nombres transfinis. Nous l'étudierons à partir
du \olref[ordinals][]{chap}. Mais, si l'on veut simplement
\emph{faire de l'arithmétique}, les deux approches conviendraient
tout aussi bien. Si quelqu'un vous dit que les nombres naturels
\emph{sont} des ensembles, la question évidente est donc :
\emph{Quels ensembles sont-ils ?}

Cette question précise a été rendue célèbre par \citet{Benacerraf1965}.
Il faut cependant souligner qu'elle n'est que l'exemple le plus connu
d'un phénomène que nous avons déjà rencontré bien des fois. Voici
l'idée fondamentale : la théorie des ensembles nous donne un moyen
de \emph{simuler} plusieurs sortes d'entités «~intuitives~» : les
réels, les rationnels, les entiers relatifs et les naturels, certes,
mais aussi les couples ordonnés, les fonctions et les relations.
Toutefois, elle ne nous impose jamais une manière \emph{unique} de
les simuler. Il existe \emph{toujours} d'autres choix qui auraient,
tout simplement, aussi bien rempli ce rôle.

\end{document}
